\section{The limiting Hessian at a unique LP optimum}
\label{sec:lp-limits}

For a unique polytope optimum the diameter law gives bounded
conditioning for every fixed barrier. The logarithmic barrier has a
stronger conclusion: a positive-definite scaled matrix limit, including
at a degenerate vertex. We present it as a classical corollary of LP
central-path endpoint theory, not as a new convergence theorem.
\citet[Theorem 3.3]{adler1991} identify the dual analytic-center limit;
their Theorem 5.1 and subsequent discussion explicitly cover the
full-column-rank active matrix at a unique optimum. Weighted path
limits and analytic extension were further developed by
\citet{guler1994,halicka1999}. The compression below makes the
consequence for the equality-reduced Hessian explicit.

Let
\begin{equation}
 P=\{x\in\R^n:Ax=b,\ x\geq0\},\qquad
 A\in\R^{m\times n},\quad\operatorname{rank}A=m<n,
 \label{eq:lp-standard}
\end{equation}
be compact and contain a strictly positive point. Suppose $c^\top x$
has a unique optimum $x^*$. Use $F_{\log}(x)=-\sum_i\log x_i$ and an
orthonormal basis matrix $W\in\R^{n\times(n-m)}$ of $\ker A$.
The centrality equations supply positive slacks $s(\mu)$ and
multipliers $y(\mu)$ with
\begin{equation}
 Ax(\mu)=b,\qquad A^\top y(\mu)+s(\mu)=c,
 \qquad x_i(\mu)s_i(\mu)=\mu.
 \label{eq:lp-centrality}
\end{equation}
In particular strict dual feasibility holds. Define
$B=\{i:x_i^*>0\}$ and $N=\{1,\ldots,n\}\setminus B$.
Classical endpoint theory gives
\begin{equation}
 s(\mu)\longrightarrow s^a,\qquad s_B^a=0,\quad s_N^a>0.
 \label{eq:lp-slack-limit}
\end{equation}
Here $s^a$ is the analytic center of the dual optimal slack face.
The strict inequality on $N$ is the strict complementarity of the
centered LP optimum; it does not assert primal nondegeneracy.

\begin{proposition}[A classical endpoint corollary]
\label{prop:lp-limit}
For the data above,
\begin{equation}
 \mu^2 H_{\log}(x(\mu))
 \longrightarrow L:=W^\top\Diag((s^a)^2)W\succ0,
 \label{eq:lp-hessian-limit}
\end{equation}
where matrices represent operators in the basis $W$. Therefore
\begin{equation}
 \mu^2\lambda_j(H_{\log}(x(\mu)))\longrightarrow\lambda_j(L)
 \quad(1\leq j\leq n-m),\qquad
 \kappa_{\log}(x(\mu))\longrightarrow\kappa(L)<\infty.
 \label{eq:lp-condition-limit}
\end{equation}
The conclusion holds when $|B|<m$, the degenerate-vertex case.
\end{proposition}

\begin{proof}
Complementarity yields the exact identity
\[
 \mu^2W^\top\Diag(x(\mu)^{-2})W
 =W^\top\Diag(s(\mu)^2)W.
\]
The right side converges in operator norm by
\eqref{eq:lp-slack-limit}. To prove that its limit is positive
definite, first observe that $A_B$ has full column rank. Otherwise a
nonzero $v\in\ker A_B$ gives feasible points
$(x_B^*\pm tv,0_N)$ for sufficiently small $t>0$. Since
$c_B=A_B^\top y^a$ by complementary slackness, both points have the
optimal objective value, contradicting uniqueness.

If $z\in\ker A$ has
$z^\top\Diag((s^a)^2)z=0$, positivity of $s_N^a$ gives $z_N=0$.
Then $A_Bz_B=0$ and full column rank imply $z=0$. Hence the quadratic
form is positive definite on $\ker A$. Continuity of ordered
eigenvalues and of the condition number on positive-definite matrices
proves \eqref{eq:lp-condition-limit}.
\end{proof}

\subsection{An active-set formula and quantitative constants}

Write $a_i$ for column $i$ of $A$. The limiting slack can be computed
from the active set without following the central path. It is
$s^a=c-A^\top y^a$, where
\begin{equation}
 \begin{aligned}
 y^a=\arg\max_y\quad &\sum_{i\in N}\log(c_i-a_i^\top y)\\
 \text{subject to}\quad&A_B^\top y=c_B,\qquad
 c_N-A_N^\top y>0.
 \end{aligned}
 \label{eq:dual-face-center}
\end{equation}
Every feasible multiplier is optimal because $b=A_Bx_B^*$.
For clarity, existence and uniqueness of this maximizer also follow
directly once strict complementarity is known. If $z>0$ is a fixed
primal feasible point and $s$ is dual optimal, then
$z^\top s=c^\top z-\OPT$ bounds every coordinate of $s$.
Full row rank makes $y\mapsto A^\top y$ injective, so the closed dual
optimal face is bounded. The logarithmic objective diverges to
$-\infty$ where an $N$-slack vanishes. On the tangent space
$\ker A_B^\top$, its negative Hessian is
\[
 A_N\Diag(s_N^{-2})A_N^\top.
\]
It is positive definite on that tangent space: if both $A_N^\top h$
and $A_B^\top h$ vanish, full row rank gives $h=0$. Thus the
maximizer is unique, with the zero-dimensional case interpreted as a
single feasible multiplier.

Let $R_N$ restrict a vector to its $N$-coordinates, and define
\[
 \sigma_N=\sigma_{\min}(R_NW)>0,\qquad
 \eta_N=\norm{R_NW}_2\leq1,\qquad
 s_{\min}=\min_{i\in N}s_i^a,\quad
 s_{\max}=\max_{i\in N}s_i^a.
\]
The same uniqueness argument proves injectivity of $R_NW$.
The bounds
\begin{equation}
 s_{\min}^2(R_NW)^\top(R_NW)
 \preceq L\preceq
 s_{\max}^2(R_NW)^\top(R_NW)
 \label{eq:lp-active-loewner}
\end{equation}
give
\begin{align}
 s_{\min}^2\sigma_N^2&\leq\lambda_{\min}(L)
       \leq s_{\max}^2\sigma_N^2,\label{eq:lp-min-constant}\\
 s_{\min}^2\eta_N^2&\leq\lambda_{\max}(L)
       \leq s_{\max}^2\eta_N^2,\label{eq:lp-max-constant}\\
 \left(\frac{s_{\min}}{s_{\max}}\right)^2
 \left(\frac{\eta_N}{\sigma_N}\right)^2
 &\leq\kappa(L)\leq
 \left(\frac{s_{\max}}{s_{\min}}\right)^2
 \left(\frac{\eta_N}{\sigma_N}\right)^2.
 \label{eq:lp-limit-bounds}
\end{align}
These constants separate reduced-cost spread from the angle of
$\ker A$ to the coordinate subspace supported on $B$. Neither factor
is controlled by the statement that the optimum is unique.
The simplex in \cref{prop:simplex-plateau} already shows unbounded
limiting condition numbers within compact instances of fixed dimension.

\begin{example}[A compact degenerate vertex]
\label{ex:degenerate-lp}
Take
\begin{equation}
 A=\begin{pmatrix}1&1&1&1\\0&1&3/2&3\end{pmatrix},\qquad
 b=\binom{1}{3/2},\qquad c=(1,1,0,1)^\top.
 \label{eq:degenerate-lp-data}
\end{equation}
The simplex equality makes $P$ compact. The point
$(2,3,3,3)^\top/11$ is strictly feasible, and $y=(-1,0)^\top$
is strictly dual feasible. The unique optimum is $e_3$, so
$|B|=1<m=2$. Every dual optimum is $y=(-3t/2,t)^\top$, with
\[
 (s_1,s_2,s_4)=(1+3t/2,1+t/2,1-3t/2),\qquad -2/3<t<2/3
\]
in its relative interior. Differentiating the sum of the three logarithms
gives $27t^2+36t-4=0$. Only one root lies in the displayed interval:
\begin{equation}
 t_a=\frac{-6+4\sqrt3}{9}.
 \label{eq:degenerate-lp-center}
\end{equation}
Strict concavity makes it the dual analytic center. For an exact
calculation without choosing an orthonormal nullspace basis, write
\[
 z=Tu,\qquad
 T=\begin{pmatrix}-1/3&1\\1&0\\-2/3&-2\\0&1\end{pmatrix},
 \qquad T^\top T=\begin{pmatrix}14/9&1\\1&6\end{pmatrix}.
\]
The columns of $T$ span $\ker A$. The eigenvalues of $L$ are the
generalized eigenvalues of
\[
 T^\top\Diag((s^a)^2)T
 =\begin{pmatrix}
 (s_1^a)^2/9+(s_2^a)^2&-(s_1^a)^2/3\\
 -(s_1^a)^2/3&(s_1^a)^2+(s_4^a)^2
 \end{pmatrix}
\]
against $T^\top T$. Numerically these are approximately
$0.24888160$ and $1.14289062$, giving $\kappa(L)\approx4.59210571$.
The exact matrix formulas establish the finite limit; the decimals only
identify its size.
\end{example}

The same endpoint proof applies to a strictly primal--dual feasible LP
with an unbounded feasible region, whenever the classical central path
has the stated centered limits. For example, for $M\geq1$,
\[
 A=\begin{pmatrix}1&0&0&0\\0&1&-1&0\end{pmatrix},\qquad
 b=\binom10,\qquad c=(0,1,1,M)^\top,
\]
the unique optimum $e_1$ is degenerate and
$x(\mu)=(1,\mu,\mu,\mu/M)^\top$,
$s(\mu)=(\mu,1,1,M)^\top$ solve the centrality equations exactly.
In the orthonormal tangent basis
$(0,1,1,0)^\top/\sqrt2$, $(0,0,0,1)^\top$, one has
$\mu^2H_{\log}=\Diag(1,M^2)$. Every positive objective sublevel is
compact, so \cref{prop:localization} also applies on its central tail.
This example shows directly that even at a degenerate vertex the finite
limit can be arbitrarily ill-conditioned.
